\section{离散对冲与模型失效}
\label{ch:lower-derivatives-hedging}

\begin{itemize}
  % Generated from curriculum/manifest.json; do not edit by hand.
\item 直接先修：下册《定价与校准：从闭式、树和 PDE 到参数化曲面》。

\end{itemize}

连续复制允许不断调整 Delta，实际交易只能在有限时点成交。下面把股票与现金逐笔相加，说明对冲误差如何产生，也说明手续费为什么会积累。


\MFQLead{Greeks 是局部敏感度，不是损益保证}

在 Black--Scholes 模型中
\[
 \Delta=\frac{\partial C}{\partial S}=\Phi(d_1),\qquad
 \Gamma=\frac{\partial^2C}{\partial S^2},\qquad
 \mathrm{Vega}=\frac{\partial C}{\partial\sigma}.
\]
解析 Greek 应与中心差分在一组逐步缩小但不过度受舍入支配的步长上比较。差分一致只检查导数实现；若价格模型漏掉股息、跳跃或曲面动态，Greek 仍可能在经济上错误。

% Generated by tools/build_figures.py; edit figures/figure-specs.json instead.
\MFQFigure{tex/figures/generated/lower-dv-greeks.pdf}{Delta、Gamma 与 Vega 描述不同局部扰动方向；同一头寸的风险会随标的价格变化而重新分配。}{lower-dv-greeks}


教学实验以现货步长 $1,0.5,0.25,0.125$ 检查 Delta 与 Gamma，并按比例缩小波动率步长检查 Vega。由最粗到最细步长，三者绝对差分别从
$6.51\times10^{-5}$、$3.10\times10^{-6}$、$6.53\times10^{-4}$ 收敛到
$1.02\times10^{-6}$、$4.84\times10^{-8}$、$1.02\times10^{-5}$；notebook 同时保存并绘制全部四个步长。步长继续缩小时舍入误差可能反弹，因此不以单个巧合步长充当证明。

小变动的 P\&L 展开为
\[
 \Delta C\approx \Delta\,\Delta S+\tfrac12\Gamma(\Delta S)^2
 +\mathrm{Vega}\,\Delta\sigma+\Theta\,\Delta t.
\]
Delta 中性只消除一阶现货项。Gamma、波动率、时间衰减、离散调仓和交易成本仍进入损益。

若实际实现波动率为 $\sigma_{\mathrm{real}}$，而期权按 $\sigma_{\mathrm{impl}}$ 定价，对于持有期权并做空 Delta 股票的方向，连续小步近似下 P\&L 的核心项与
\[
 \frac12\Gamma S^2\left(\sigma_{\mathrm{real}}^2-\sigma_{\mathrm{impl}}^2\right)\Delta t
\]
相关。它解释了波动率观点如何进入对冲损益，但仍忽略跳跃、曲面重标定、交易成本和 Greek 离散误差，不能当作逐路径恒等式。

\MFQLead{自融资离散台账}

路径仍由风险中性 Brownian 运动 $W_t$ 驱动；实际对冲评估则必须说明使用哪个测度和真实世界路径模型。

期初卖出一份期权，买入 $\Delta_0$ 股，其余进入现金账户。每个调仓时点先让现金按无风险利率增长，再用当前价格买卖 $\Delta_{t_i}-\Delta_{t_{i-1}}$ 股，并扣除
\[
 \mathrm{cost}_{t_i}=c\,S_{t_i}\left|\Delta_{t_i}-\Delta_{t_{i-1}}\right|.
\]
到期组合价值减支付收益得到复制误差。必须保存旧 Delta、交易量、现货价格、现金、成本和到期 payoff；只保存最终误差无法定位符号、融资或时点错误。

以一步路径为手算基准时，期初成本为 $c|\Delta_0|S_0$，到期前不再调仓，故成本拖累应等于该成本按现金账户增长后的值。仓库保留这条独立手算测试，用来捕获“成本没有融资”“到期重复调仓”或“成本符号加反”等错误。多路径实验是在这个小计算记录通过以后才有意义。

透明实现逐路径执行该计算记录，第二个独立向量实现以 SciPy \texttt{ndtr} 计算批量 Delta，并用 NumPy 同时更新 512 条路径。两者最大逐路径误差差约为 $5.7\times10^{-14}$；它们共享模型假设，不能互相证明 GBM 合理。

\MFQLead{从一条路径升级为分布}

单路径误差可能恰好很小或很大。固定 512 条路径、24 个时间步后，至少报告
\[
 \mathrm{bias}=N^{-1}\sum_j e_j,\qquad
 \mathrm{RMSE}=\sqrt{N^{-1}\sum_j e_j^2},
\]
以及 5\%、50\%、95\% 分位数。教学实验无成本 bias/RMSE 为 $(-0.086348,1.416360)$；加入比例成本后为 $(-0.192103,1.434142)$，误差分位数为 $(-2.485457,-0.052646,1.960850)$。平均名义成本为 $0.104415$，95\% 分位为 $0.149720$。

% Generated by tools/build_figures.py; edit figures/figure-specs.json instead.
\MFQFigure{tex/figures/generated/lower-dv-hedge-error.pdf}{离散对冲误差在不同路径上形成一个分布；在本例模型中，提高调仓频率减少离散误差，却可能增加交易成本。}{lower-dv-hedge-error}


这些数值只描述固定随机种子和参数。正式实验还应改变调仓频率、成本、跳跃强度、实现波动率、隐含波动率与曲面动态，报告误差的抽样不确定性。提高调仓频率通常降低离散误差，却可能增加成本；不存在脱离摩擦的“越频繁越好”。

同一规则在 12、24、52 次调仓下的成本后 RMSE 为 $2.040789$、$1.434142$、$0.919696$，平均名义成本为 $0.081246$、$0.104415$、$0.139579$。该组实验同时展示离散误差下降与成本上升，不能从中挑一个维度宣称“更优”。

bias 与 RMSE 本身也有抽样误差。可以对路径级误差 bootstrap，或用多组独立 seed 重复整个实验并报告区间。若路径共享市场因子或波动率状态，简单把每条路径当独立观察会低估不确定性；分组重采样单位应与共同冲击结构一致。

\MFQTransition{跳跃与离散调仓}
跳跃、随机波动率和延迟使实际路径偏离连续复制条件。对固定策略比较不同调仓频率与成本率，应同时观察误差分布和总交易成本；路径多只能减小给定模型内的抽样误差。

\MFQTransition{Greek 的单位与组合聚合}

Delta 是价格对标的的一阶导，Gamma 是二阶曲率，Vega 是对波动率“绝对变化 1”还是“一个波动率点”的敏感度必须注明。Theta 的日历日/交易日单位、Rho 的利率单位也常被混淆。组合 Greek 先按合约乘数、数量和币种换算后再相加。

有限差分检查采用中心差分
\[
\Delta\approx\frac{V(S+h)-V(S-h)}{2h},\qquad
\Gamma\approx\frac{V(S+h)-2V(S)+V(S-h)}{h^2}.
\]
$h$ 太大截断误差高，太小则浮点消减严重；应画 $h$ 的误差曲线，而不是只挑一个能对上的步长。

\MFQTransition{离散 Delta 对冲误差的来源}

持有期权、做空 Delta 股票的损益可局部近似分解为
\[
\Delta\Pi\approx
\frac12\Gamma S^2(\sigma_{\rm realized}^2-\sigma_{\rm hedge}^2)\Delta t
-\text{cost}+\text{jump/model residual}.
\]
后文卖出期权并持有复制组合的误差方向相反。这是局部近似，揭示 Gamma 暴露如何把实现方差与对冲波动率差转成损益。更频繁调仓降低离散误差，却增加点差和冲击，因此存在成本依赖的最优频率，而不是“越频繁越好”。

跳跃时 Delta 无法提前复制不连续变动；随机波动率引入 Vega 风险；曲面移动还涉及 sticky-strike、sticky-delta 等约定。对冲报告应按这些情景分组，而非只给总体 RMSE。

\MFQTransition{从模型价格到市场报价}

做市价格还要覆盖库存、资本、模型不确定性和流动性。市场 bid/ask 不是同一个“真波动率”的噪声观察；用中价校准后按中价计算对冲收益会忽略入场和退出价差。

曲面随时间更新时，参数跳动可能来自数据噪声、报价稀疏或优化局部极值。应记录报价过滤、权重、初值、约束、收敛状态和参数时间序列。无法校准时应保留失败状态，不用昨日参数静默填充。

\MFQTransition{对冲实验的统计设计}

多路径实验至少报告 bias、RMSE、中位数、尾部分位、成本分布和最大资金需求。比较两个对冲策略时使用共同随机数，让路径噪声配对抵消。模型和策略调参使用开发情景，最终压力情景只打开一次。

真实历史回放中只有一条市场路径，可按滚动起点形成多个重叠对冲 episode，但标准误要处理重叠。历史回放不能评价订单对市场的反事实影响，成交模型仍需单独压力测试。

\begin{diagnostic}
对冲误差小不等于定价模型正确：高频调仓和低成本假设可以掩盖模型错设。必须同时报告交易次数、成本、未成交、资本占用和压力期表现。
\end{diagnostic}

\section{本路线的综合项目}

以同一合约和参数贯通复制、定价、校准与离散对冲。路线 Capstone 的综合结果应给出价格与 Greeks 的独立核对、隐波或曲面的适用范围，以及多路径对冲误差、成本和压力情景分布；单条漂亮路径不能替代这些信息。

\subsection{两次持仓与现金更新}

设卖出期权收到价格 10，$S_0=100$，初始 Delta 为 $0.6$，利率暂为零、成本率为 $c=0.001$。买股花费 60，初始费用 0.06，现金为 $B_0=10-60-0.06=-50.06$。下一时点价格 110，新的 Delta 为 $0.8$；加买 0.2 股花费 22，费用 0.022，所以 $B_1=-72.082$。

到期 $S_T=105,K=100$，若将股票按中价估值而不另收平仓费，复制组合价值为 $0.8(105)-72.082=11.918$，支付为 5，误差为 $6.918$。这些 Delta 是教学输入，并非声称与某一波动率模型匹配。该例完整说明现金流符号，实际模型 Delta 会在 notebook 中由价格公式计算。

一般递推为 $B_i=B_{i-1}e^{r\Delta t_i}-S_i(\Delta_i-\Delta_{i-1})-cS_i|\Delta_i-\Delta_{i-1}|$。到期不再重新估计一个用于未来持有的 Delta，而是计算 $B_{m-1}e^{r\Delta t_m}+\Delta_{m-1}S_T-g(S_T)$；若要求实际平仓，再单独扣平仓费用。notebook 比较多条路径和不同调仓频率，保留成本前后误差的配对关系。

\section{分层练习}
\begin{enumerate}
  \item \textbf{口述概念：}树、PDE 和 Monte Carlo 各自主要有哪些误差？
  \item \textbf{口述概念：}为什么逐点隐波反演不是参数化曲面校准？
  \item \textbf{笔试推导：}写出隐式 PDE 三对角系统的三个系数。
  \item \textbf{笔试推导：}推导非等距执行价网格的斜率凸性条件。
  \item \textbf{数值编程：}分别加密空间、时间和截断边界，制作 PDE 误差表。
  \item \textbf{数值编程：}把参数化总方差系数恢复与价格残差分开报告。
  \item \textbf{研究判断：}何时固定执行价的日历单调性不宜直接使用？
  \item \textbf{研究判断：}财政部 par yield 快照还缺哪些输入才能成为期权贴现曲线？
\end{enumerate}

\MFQLead{Greeks、离散对冲与误差分布题组}
\begin{enumerate}
  \item \textbf{口述概念：}Delta 中性以后还剩哪些风险？
  \item \textbf{口述概念：}为什么单路径复制误差不能评价对冲质量？
  \item \textbf{笔试推导：}写出包含融资与比例成本的离散自融资递推。
  \item \textbf{笔试推导：}说明 bias、RMSE 和分位数回答的不同问题。
  \item \textbf{数值编程：}把调仓步数改为 12、24、52，并比较误差与成本。
  \item \textbf{数值编程：}加入一处价格跳跃，比较成本前后误差分布。
  \item \textbf{研究判断：}审计只展示一条“完美复制”路径的报告。
  \item \textbf{开放 Capstone：}结合本章模型，解释预测、成交和成本如何共同决定结果，并说明所用假设。
\end{enumerate}

\subsection*{基础衔接：独立计算}
在上述例子中增加到期平仓费用，误差减少多少？

